\documentclass[12pt, titlepage]{article}

\usepackage{booktabs}
\usepackage{tabularx}
\usepackage{float}
\usepackage{hyperref}
\usepackage{longtable}
\usepackage{array}
\usepackage{enumitem}
\usepackage{geometry}
\usepackage{ragged2e}
\usepackage{caption}

\geometry{margin=1in}
\setlength{\parskip}{0.4em}
\setlength{\parindent}{0pt}

\hypersetup{
    colorlinks,
    citecolor=black,
    filecolor=black,
    linkcolor=red,
    urlcolor=blue
}

\input{../../Comments}
%% Common Parts

\newcommand{\progname}{Project Proxi} % PUT YOUR PROGRAM NAME HERE
\newcommand{\authname}{Team \#23, Project Proxi
\\ Savinay Chhabra
\\ Amanbeer Singh Minhas
\\ Gourob Podder
\\ Ajay Singh Grewal} % AUTHOR NAMES                  

\usepackage{hyperref}
    \hypersetup{colorlinks=true, linkcolor=blue, citecolor=blue, filecolor=blue,
                urlcolor=blue, unicode=false}
    \urlstyle{same}
                                


\begin{document}

\hypersetup{pageanchor=false}

\title{Usability Testing Report: Project Proxi}
\author{\authname}
\date{\today}

\maketitle

\pagenumbering{roman}

\section{Revision History}

\begin{tabularx}{\textwidth}{p{3cm}p{2cm}X}
\toprule
{\bf Date} & {\bf Version} & {\bf Notes}\\
\midrule
March 20th, 2026 & 1.0 &
Initial draft created with testing scope, plan, and result
placeholders.\\
March 23rd, 2026 & 1.1 &
Added participant profiles, task scenarios, questionnaires,
and data analysis plan.\\
March 26th, 2026 & 1.2 &
Added quantitative and qualitative findings, themes,
severity ratings, and recommendations.\\
March 27th, 2026 & 2.0 &
Added exploratory free-use sessions, feature-discovery
analysis, and updated appendices.\\
\bottomrule
\end{tabularx}

\newpage

\tableofcontents

\newpage

\hypersetup{pageanchor=true}
\pagenumbering{arabic}

\section{Introduction}

Project Proxi is an AI-powered desktop assistant designed to help users
complete computer tasks through natural voice and text interaction.
Proxi was developed to improve accessibility, efficiency, and ease of
use for seniors, users with disabilities, and power users who want a
faster workflow for common digital tasks.

The purpose of this report is to present the usability testing
conducted for Proxi. This report follows a systematic, rigorous, and
quantified evaluation process. The study includes a clear testing plan,
realistic task scenarios, structured pre-test and post-test questions,
post-task rating instruments, observed task metrics, and qualitative
analysis based on participant feedback.

The main objective of this usability study was to determine whether
Proxi is understandable, useful, and efficient when users attempt
realistic daily workflows such as reading and sending emails, creating
and managing calendar events, requesting weather information, updating
notes, and interacting with the system using speech. Because Proxi was
designed to support both accessibility-focused and productivity-focused
usage, the study evaluated not only task success but also confidence,
clarity, comfort, learnability, and willingness to use the system in
real life.

\section{Executive Summary}

This usability study examined how well Project Proxi supports realistic
computer tasks through a high-level conversational interface. Four
participants were recruited to represent different intended user groups:
a power user, an elderly user, a voice-first professional user, and a
general daily-task user. Participants completed scenario-based tasks
covering email, calendar, notes, search, weather, and accessibility
features. In addition to the structured task study, two participants
were also given a short exploratory session with open access to Proxi.
During this session, they were encouraged to interact with the system
freely and try whatever features felt natural to them. This helped
capture spontaneous usage patterns, feature discovery, and unprompted
usability issues that may not appear in scripted tasks.

The results showed that Proxi performed strongly in workflows involving
email drafting, calendar management, weather retrieval, and note
creation. Participants especially appreciated the convenience of using
one assistant to complete multi-step actions without switching between
many applications. They also responded positively to the voice-first
interaction style and to the system's potential to support users who
find traditional interfaces frustrating.

The most significant usability issue observed was command
discoverability. Participants often knew what they wanted to do, but
were not always sure how to phrase the request. Some follow-up requests
were also too ambiguous for the system to resolve smoothly. In addition,
Punjabi support and other accessibility-oriented features were seen as
highly valuable, but participants wanted such capabilities to be made
more visible earlier in the interaction.

Overall, the findings suggest that Proxi is highly promising as both an
accessibility tool and a productivity assistant. The report concludes
with design recommendations focusing on onboarding, example prompts,
multilingual support, contextual clarification, and better visibility
of supported commands.

\section{Overview}

Project Proxi is intended to reduce the effort required to perform
common computer tasks. Instead of relying only on a keyboard, mouse, or
multiple application interfaces, users can ask Proxi to complete tasks
through text or speech. These tasks include reading unread emails,
drafting replies, showing calendars, creating events, updating notes,
summarizing information, and answering practical day-to-day questions.

Usability is a central concern for this project because the target users
include seniors and accessibility-focused users who may struggle with
small controls, multi-step actions, and complex interface navigation.
At the same time, power users expect the system to be fast, flexible,
and efficient enough to be worth adopting into their daily workflow.

This report therefore focuses on whether Proxi supports the right kinds
of interaction for its intended audience and whether the current design
needs refinement before final delivery.

\section{Usability Testing Goals and Success Criteria}

This study was designed to evaluate Proxi in a systematic, rigorous,
and quantified manner consistent with the expectations of the HCI
course. The testing plan focused on whether representative users could
complete realistic tasks, whether the interaction style was easy to
learn, and whether the system felt useful enough for real use.

\subsection{Usability Testing Goals}

The main goals of the usability study were as follows:
\begin{enumerate}[label=\arabic*.]
    \item Evaluate learnability by determining whether first-time users
    can understand how to begin interacting with Proxi after a short
    orientation.
    \item Evaluate task success by measuring whether users can complete
    realistic workflows involving email, calendar, search, notes, and
    weather.
    \item Evaluate accessibility by observing whether an elderly user
    and a lower-English-fluency user can interact with the system with
    confidence and low frustration.
    \item Evaluate efficiency by observing whether Proxi reduces the
    need to switch between applications for power and professional
    users.
    \item Evaluate exploratory use by observing what users try when they
    are given open access to Proxi without being told exactly what to
    do.
    \item Evaluate clarity of system response by determining whether
    users understand what Proxi is doing and what happened after each
    command.
    \item Evaluate confidence and trust by measuring whether users feel
    that Proxi understood their requests and behaved predictably.
    \item Evaluate satisfaction by measuring whether participants would
    use Proxi again in real life.
\end{enumerate}

\subsection{Success Criteria}

To make the evaluation more rigorous, the team defined the following
success criteria before reviewing the study results:

\begin{itemize}
    \item unaided task completion rate of at least 80\% across all core
    tasks.
    \item average overall satisfaction score of at least 4.0 out of 5.
    \item average confidence score of at least 4.0 out of 5 for whether
    Proxi understood user intent.
    \item average assists per task of no more than one assist for every
    two tasks.
    \item successful completion of at least one key task for each target
    user type represented in the study.
    \item clear qualitative evidence that participants understood the
    value of Proxi for either accessibility, efficiency, or both.
\end{itemize}

\subsection{Administration of Measures}

The same procedure was used for each participant. One moderator guided
the session and one note-taker recorded times, attempts, assists,
confusion points, and direct quotes. After each task, participants
completed the same post-task questionnaire using a 1 to 5 Likert scale
before moving to the next task. At the end of the session, each
participant completed the same post-test interview.

\subsection{Justification of Study Size}

This was a formative usability study intended to identify major
usability issues across a small set of representative user types rather
than to produce statistically generalizable results. The participant set
was intentionally selected to cover distinct usage contexts that matched
the core Proxi user groups, including an elderly user, a power user, a
voice-first professional user, and a general daily-task user.

\section{Research Questions}

The usability study was designed to answer the following questions:

\begin{enumerate}[label=\arabic*.]
    \item Can users complete realistic workflows involving email,
    calendar, search, weather, and notes?
    \item Does Proxi support both accessibility-driven and
    productivity-driven usage patterns?
    \item Are users satisfied with the overall conversation flow?
    \item What barriers do users encounter while speaking to or typing
    commands into Proxi?
    \item Do users understand what Proxi can do without relying too
    heavily on help prompts?
    \item Are the responses clear, actionable, and trustworthy?
    \item What improvements would make Proxi easier to use in everyday
    settings?
    \item What features do users discover and attempt on their own when
    they are given open-ended access to Proxi?
\end{enumerate}

\section{Participant Profiles}

Participants were selected to reflect the intended user space of Proxi.
Each participant represented a distinct use context that matched one or
more project personas and user groups.

\subsection{Participant 1: Raj - Power User}

Raj is a tech-savvy user who regularly works with research, note-taking,
and email. He often repeats multi-step workflows and wants a system that
reduces application switching. His main goal was to test whether Proxi
could speed up searching, summarizing, note updates, and email actions.
Raj also participated in the exploratory free-use portion of the study
because he was likely to test a broader range of productivity features
without needing heavy prompting.

\subsection{Participant 2: Tony - Elderly User}

Tony is an older user who prefers a light theme and simpler interaction.
He speaks less English and is more comfortable with direct spoken
interaction than with typing. His scenario emphasized accessibility,
theme preference, ease of understanding, and multilingual support.

\subsection{Participant 3: Jassi - Voice Productivity User}

This participant works in real estate and frequently needs to schedule
meetings and send follow-up messages while moving between appointments.
His scenario focused on the usefulness of voice-based calendar and email
workflows under time pressure.

\subsection{Participant 4: Mikayla - Daily-Task User}

The fourth participant represented a broader everyday use case. This
scenario focused on practical daily information tasks, including
weather, reminders, and lightweight note creation. Mikayla also
participated in the exploratory free-use portion of the study to help
evaluate how a general user approaches Proxi when given open-ended
access rather than specific scripted tasks.

\subsection{Recruitment Criteria}

Participants were included if they matched at least one of the
following:
\begin{itemize}
    \item older adult or accessibility-oriented user.
    \item regular user of email, calendar, or note workflows.
    \item someone likely to benefit from voice interaction.
    \item willing participant in a moderated usability study.
\end{itemize}

Participants were excluded if they were directly involved with Proxi's
development or were already overly familiar with the exact command set.

\section{Methodology}

This study used a moderated task-based usability testing approach. Each
session followed a consistent structure:
\begin{enumerate}[label=\arabic*.]
    \item introduction and consent.
    \item pre-test questionnaire.
    \item short orientation.
    \item warm-up task.
    \item scenario-based tasks.
    \item post-task questionnaire after each task.
    \item post-test interview.
    \item final notes and debrief.
\end{enumerate}

This method was chosen because Proxi is an interactive system whose
value depends on real workflows rather than isolated feature checks. A
moderated study allowed the team to record both measurable task results
and richer participant reactions such as hesitation, delight,
confusion, and suggestions for improvement.

\subsection{Session Length}

Each usability session lasted approximately 25 to 35 minutes.

\subsection{Data Collected}

For each participant and task, the following data was recorded:
\begin{itemize}
    \item task completion status.
    \item time to complete the task.
    \item number of attempts.
    \item number of retries.
    \item number of moderator assists.
    \item post-task ratings.
    \item direct participant quotes.
    \item notable confusion points.
    \item final comments and recommendations.
    \item features attempted during exploratory free use.
    \item number of spontaneous commands not tied to scripted tasks.
    \item order of feature discovery during exploratory use.
\end{itemize}

\subsection{Exploratory Free-Use Session}

In addition to the structured task sessions, two participants completed
a short exploratory free-use session. These participants were given open
access to Proxi for approximately 8 to 10 minutes after their main task
session. They were told that they could ``play around'' with the system
and try any features that felt natural, useful, or interesting to them.
The moderator did not provide specific tasks during this part of the
study unless the participant asked a clarifying question about system
availability.

The purpose of this exploratory session was to observe spontaneous
feature discovery, natural command phrasing, and unprompted breakdowns.
This portion of the study was especially useful for understanding what
participants believed Proxi could do without being guided by scripted
prompts.

During the exploratory session, the note-taker recorded:
\begin{itemize}
    \item which features participants chose to try first,
    \item how often they switched between voice and text input,
    \item whether they used the help feature,
    \item whether they returned to features from the structured tasks,
    \item moments of surprise, delight, hesitation, or confusion,
    \item direct quotes describing what they expected Proxi to do.
\end{itemize}

\section{Testing Environment and Setup}

Testing took place in a quiet indoor environment using a laptop running
a working Proxi build with microphone input enabled. The moderator sat
beside the participant and guided the session. A note-taker recorded
observations, quotes, retries, assists, and time measurements.

The hardware and materials used during testing were:
\begin{itemize}
    \item one laptop running Proxi.
    \item one microphone for speech input.
    \item keyboard for optional text interaction.
    \item printed task packet.
    \item printed rating sheets.
    \item observation sheets for note-taking.
\end{itemize}

The environment was controlled to reduce interruptions and maintain
similar testing conditions across all participants.

\section{Pre-Test Questionnaire}

Before beginning the tasks, participants were asked the following
questions.

\subsection{General Information}

\begin{enumerate}[label=\arabic*.]
    \item What is your age range?
    \item Which role best describes you: student, working professional,
    retiree, or other?
    \item How comfortable are you with computers in general?
    \item How often do you use voice assistants such as Siri, Alexa,
    Google Assistant, or dictation tools?
    \item Do you usually prefer typing, speaking, or a mix of both?
\end{enumerate}

\subsection{Current Workflow}

\begin{enumerate}[label=\arabic*.]
    \setcounter{enumi}{5}
    \item What computer tasks do you do most often in a normal week?
    \item Which of these do you use most: email, calendar, notes, file
    reading, web search, reminders?
    \item What usually feels slow or frustrating in your current
    workflow?
    \item Do you often repeat the same steps across different apps?
    \item Have you ever wished you could just tell the computer what to
    do?
\end{enumerate}

\subsection{Expectations}

\begin{enumerate}[label=\arabic*.]
    \setcounter{enumi}{10}
    \item Before using Proxi, what tasks would you expect a tool like
    this to help with?
    \item What would make a system like this feel genuinely useful to
    you?
    \item What would make you not trust it?
\end{enumerate}

\subsection{Accessibility and Comfort}

\begin{enumerate}[label=\arabic*.]
    \setcounter{enumi}{13}
    \item Do you ever find menus, small buttons, or multi-step actions
    frustrating?
    \item Would voice input be useful in your daily life? Why or why
    not?
    \item Would language flexibility, such as Punjabi support, make a
    difference for you or someone you know?
\end{enumerate}

\section{Training and Warm-Up Procedure}

Before the main tasks, each participant received a short orientation.
The moderator explained that:
\begin{itemize}
    \item Proxi supports both voice and text input.
    \item The participant was not being tested.
    \item There was no single perfect phrasing.
    \item Confusion or frustration should be described openly.
    \item The participant could skip a task if needed.
\end{itemize}

Each participant then completed one simple warm-up task that was not
counted in the study results. The warm-up task was either:
\begin{itemize}
    \item ask Proxi what the weather is today; or
    \item ask Proxi what it can do.
\end{itemize}

This allowed the participant to hear a response and get comfortable
before the formal task timing began.

\section{Task Scenarios}

Each participant completed realistic tasks written as scenarios rather
than as rigid commands. This was done to evaluate whether users could
naturally express their goals.

\subsection{Scenario 1: Raj - Power Workflow}

Raj is in the middle of a busy workday. He wants to do quick research,
save useful information to a note page, and handle email without opening
several different applications.

\textbf{Tasks}
\begin{enumerate}[label=\arabic*.]
    \item Search the web for a topic and summarize the result.
    \item Add the summary to a note or Notion page.
    \item Read the latest unread email.
    \item Draft a reply email based on the context.
\end{enumerate}

\subsection{Scenario 2: Tony - Elderly and Punjabi-Support Workflow}

Tony wants an easier way to use the computer. He prefers light mode,
speaks less English, and benefits from simpler spoken interaction.

\textbf{Tasks}
\begin{enumerate}[label=\arabic*.]
    \item Change the application to light theme.
    \item Ask Proxi for today's weather using speech.
    \item Ask a follow-up question in a simpler or mixed-language way.
    \item Request the answer again in an easier-to-understand form.
    \item Use the help feature if unsure what to say next.
\end{enumerate}

\subsection{Scenario 3: Realtor - On-the-Go Scheduling}

The participant is moving between appointments and wants to update the
day quickly using voice.

\textbf{Tasks}
\begin{enumerate}[label=\arabic*.]
    \item Show tomorrow's calendar.
    \item Create a client meeting.
    \item Change or delete a conflicting event.
    \item Draft an email confirming the meeting.
\end{enumerate}

\subsection{Scenario 4: Everyday Utility}

The participant wants help planning the day with quick information.

\textbf{Tasks}
\begin{enumerate}[label=\arabic*.]
    \item Search for Hamilton weather.
    \item Summarize the result.
    \item Create a short daily note or page.
    \item Add a reminder or event related to the day's plan.
\end{enumerate}

\section{Post-Task Questionnaire}

After each task, participants completed the following rating questions
using a 1 to 5 Likert scale, where 1 means strongly disagree and 5 means
strongly agree.

\begin{enumerate}[label=\arabic*.]
    \item I was able to complete this task easily.
    \item I understood what the system was doing.
    \item The system responded quickly enough for me.
    \item I knew what to say or type during this task.
    \item The system feedback was clear and helpful.
    \item I felt confident that Proxi understood me.
    \item I would use this feature in real life.
    \item I was satisfied with how this task went overall.
\end{enumerate}

Participants were then asked the following short open-ended questions:
\begin{enumerate}[label=\arabic*.]
    \item What, if anything, confused you during this task?
    \item What did you like most about this task?
    \item What would have improved this task for you?
\end{enumerate}

\section{Post-Test Interview}

At the end of the session, participants answered the following
questions:

\begin{enumerate}[label=\arabic*.]
    \item What did you like most about using Proxi?
    \item What did you dislike most about using Proxi?
    \item What did you think about the overall flow of interaction?
    \item Did any feature feel surprisingly useful?
    \item Did anything feel harder than you expected?
    \item Was there any point where you were unsure how to phrase your
    request?
    \item Did you feel like the system gave enough feedback during
    waiting or processing?
    \item Would you use Proxi in the future? Why or why not?
    \item What other things would you want Proxi to help with?
    \item Is there anything else you would like to add?
\end{enumerate}

\section{Data Analysis Plan}

Both quantitative and qualitative data were analyzed.

\subsection{Quantitative Measures}

The following metrics were recorded:
\begin{itemize}
    \item time per task.
    \item number of attempts.
    \item number of retries.
    \item number of assists.
    \item task completion rate.
    \item satisfaction score.
    \item confidence score.
    \item future-use score.
\end{itemize}

\subsection{Qualitative Measures}

The following types of qualitative evidence were collected:
\begin{itemize}
    \item direct quotes.
    \item observed hesitation.
    \item confusion about phrasing.
    \item trust and confidence comments.
    \item feature requests.
    \item emotional reactions such as relief, surprise, and frustration.
\end{itemize}

\subsection{Planned Analysis}

The team used descriptive statistics for numeric measures and thematic
analysis for open-ended responses and observations. Issues were then
grouped by frequency, severity, and impact.

\subsection{Exploratory Session Analysis}

For the exploratory session, the team used descriptive counts and
thematic analysis rather than task-completion scoring alone. Since this
portion of the study was open-ended, the analysis focused on what users
chose to do, which features they discovered without prompting, which
commands felt natural to them, and what kinds of confusion arose during
self-directed use. This allowed the team to compare structured task
performance with spontaneous real-world style interaction.

\section{Quantitative Findings}

A total of 16 formal task attempts were recorded across 4 participants.

\begin{table}[H]
\centering
\caption{Overall Quantitative Results}
\begin{tabularx}{0.75\textwidth}{l>{\RaggedLeft\arraybackslash}p{3cm}}
\toprule
Metric & Result \\
\midrule
Total participants & 4 \\
Total tasks attempted & 16 \\
Completed without help & 13 \\
Completed with help & 2 \\
Not completed & 1 \\
Unaided completion rate & 81.25\% \\
Average task time & 52 seconds \\
Average attempts per task & 1.4 \\
Average assists per task & 0.31 \\
Average satisfaction & 4.3 / 5 \\
Average confidence Proxi understood intent & 4.1 / 5 \\
Average future-use interest & 4.4 / 5 \\
\bottomrule
\end{tabularx}
\end{table}

The quantitative results suggest that Proxi supports core tasks
successfully for most users. Satisfaction and future-use scores were
especially strong, indicating that participants saw practical value in
the system even when they encountered small issues.

\section{Exploratory Session Findings}

Two participants, Raj and Mikayla, completed an additional exploratory
free-use session after the structured task portion of the study. Each
participant was given approximately 10 minutes of open access to Proxi
and was encouraged to try any features that seemed useful or
interesting.

\subsection{Exploratory Quantitative Summary}

\begin{table}[H]
\centering
\caption{Exploratory Free-Use Session Summary}
\begin{tabularx}{\textwidth}{p{3cm}ccX}
\toprule
Participant & Commands Tried & Features Used & Notes \\
\midrule
Raj & 9 & 5 &
Used search, notes, email, help, and calendar features. \\
Mikayla & 7 & 4 &
Used weather, reminders, help, and note creation features. \\
\bottomrule
\end{tabularx}
\end{table}

Across the two exploratory sessions, participants issued a total of 16
self-directed commands and used 6 distinct feature types without being
prompted by the moderator. The help feature was used 3 times across the
two sessions. Both participants returned to at least one feature they
had already used during the structured task study, but they also tried
new commands that had not been assigned to them in the formal scenario
tasks.

\subsection{Exploratory Qualitative Summary}

The exploratory sessions showed that participants quickly formed mental
models about what Proxi was useful for. Raj treated Proxi primarily as a
workflow assistant and naturally tested whether it could connect search,
notes, and email. Mikayla treated Proxi more like a practical daily
assistant and explored weather, reminders, and lightweight note-taking.

The exploratory session also revealed that users were willing to test
feature boundaries. When they were not given a script, participants
more often tried broad requests such as asking what Proxi could do,
requesting multi-step actions, or testing whether earlier context would
carry into follow-up commands. This produced richer evidence about how
the system might be used in real life.

\section{Qualitative Findings}

\subsection{Successes}

\subsubsection{Strong perceived value for real tasks}

Participants repeatedly noted that Proxi felt useful when it reduced
application switching.

\begin{quote}
``This is actually nice because I don't have to open three different
things just to do one small task.''
\end{quote}

\begin{quote}
``The email and calendar part felt the most real to me.''
\end{quote}

\begin{quote}
``For work, this would save me time every day.''
\end{quote}

\subsubsection{High value for voice-first accessibility}

Tony's session showed that voice interaction reduced tension around
computer use.

\begin{quote}
``I like that I don't have to keep finding buttons.''
\end{quote}

\begin{quote}
``Light mode is much better for me.''
\end{quote}

\begin{quote}
``Oh, I didn't know Punjabi option was also there.''
\end{quote}

\begin{quote}
``If it can understand this way too, then that's actually very
helpful.''
\end{quote}

\begin{quote}
``For older people this is easier than typing.''
\end{quote}

\subsubsection{Useful for mobile-style professionals}

The realtor scenario showed strong value for voice-first scheduling.

\begin{quote}
``This is exactly the kind of thing I would use in the car before going
inside.''
\end{quote}

\begin{quote}
``Calendar plus email together is very practical.''
\end{quote}

\begin{quote}
``I liked that I could just say it instead of going back and forth.''
\end{quote}

\subsubsection{Discoverability improved after early success}

Participants generally became more confident after one successful task.

\begin{quote}
``Once I understood how it wanted me to ask, it got much easier.''
\end{quote}

\begin{quote}
``The first task felt strange, then after that I got it.''
\end{quote}

\begin{quote}
``Now I know what kind of commands it likes.''
\end{quote}

\subsection{Exploratory Use Behaviours}

The exploratory sessions provided useful evidence about how participants
approached Proxi when they were not given explicit instructions. Raj
quickly began testing whether Proxi could chain together multiple
productivity actions. After completing the structured tasks, he tried
broader commands and commented:

\begin{quote}
``Now I want to see if it can do the full thing without me breaking it
into steps.''
\end{quote}

He also tried combining search and note-taking more naturally and said:

\begin{quote}
``This is the kind of thing I would actually ask if I were working.''
\end{quote}

Mikayla's exploratory use was more curiosity-driven. She started by
asking what else Proxi could do and then moved into daily planning
features. She commented:

\begin{quote}
``I wasn't told to do reminders here, but that was the next thing I
wanted to try.''
\end{quote}

She also reacted positively when discovering additional capability:

\begin{quote}
``It feels better when I can just explore instead of following exact
steps.''
\end{quote}

At the same time, exploratory use made command discoverability issues
more visible. Without a task prompt in front of them, both participants
were more likely to pause before speaking and more likely to use the
help feature. One participant noted:

\begin{quote}
``When I have a goal in mind, I still want examples of how to ask it.''
\end{quote}

\subsection{Challenges}

\subsubsection{Difficulty knowing how to phrase requests}

This was the most important usability issue. Participants often knew the
goal but were unsure how to ask for it.

\begin{quote}
``I know what I want, I just don't know how to ask it.''
\end{quote}

\begin{quote}
``It feels like I need to guess the right sentence.''
\end{quote}

\begin{quote}
``I wanted it to understand what I meant, not the exact wording.''
\end{quote}

\begin{quote}
``I wasn't sure whether to say summarize the page, summarize the
website, or summarize the result.''
\end{quote}

\subsubsection{Ambiguity in follow-up commands}

Some participants expected natural references like ``that'' or ``do that
for tomorrow'' to work more smoothly.

\begin{quote}
``I thought it would know what `that' meant.''
\end{quote}

\begin{quote}
``I was referring to the same thing, but I guess it needed more
detail.''
\end{quote}

\begin{quote}
``The follow-up felt natural to me, but maybe not to the system.''
\end{quote}

\subsubsection{Need for more visible examples}

Participants wanted examples before needing to ask for help.

\begin{quote}
``A little box saying `try asking this' would help a lot.''
\end{quote}

\begin{quote}
``Once I used help, it made more sense.''
\end{quote}

\begin{quote}
``I wish it showed examples without me having to ask first.''
\end{quote}

\subsubsection{Multilingual support needs better onboarding}

Punjabi support was perceived as very valuable, but participants wanted
it to be introduced more explicitly.

\begin{quote}
``If I knew earlier that Punjabi was supported, I would have tried that
first.''
\end{quote}

\begin{quote}
``Oh wow, I didn't know it could understand that too.''
\end{quote}

\begin{quote}
``This would help people who mix languages naturally.''
\end{quote}

\subsubsection{Note and page updates need more flexibility}

Raj expected note titles and page references to be handled more
forgivingly.

\begin{quote}
``I knew which page I meant, but I guess the title had to match more
closely.''
\end{quote}

\begin{quote}
``This part should be smarter about similar names.''
\end{quote}

\section{Severity Ratings}

Issues were grouped by severity using a scale inspired by standard
usability severity ratings.

\subsection{Severity 4 - Usability Catastrophe}

\textbf{Difficulty with phrasing requests.}

Participants may know the intended goal but still fail because they are
not sure what wording the system accepts. This issue affects first-time
learnability and is important to fix before broader release.

\subsection{Severity 3 - Major Usability Problem}

\textbf{Ambiguous follow-up references.}

This issue did not block all tasks, but it caused noticeable breaks in
the interaction flow and reduced confidence.

\subsection{Severity 3 - Major Usability Problem}

\textbf{Hidden multilingual capability.}

Participants strongly valued Punjabi support once discovered, but its low
visibility limited its benefit.

\subsection{Severity 2 - Minor Usability Problem}

\textbf{Theme preference discoverability.}

Participants liked the light theme once they saw it, but it was not
always obvious where to find or how to change it.

\subsection{Severity 2 - Minor Usability Problem}

\textbf{Limited visible examples before help request.}

This did not stop most tasks, but it slowed early confidence and made
first impressions weaker.

\section{Themes and Recommendations}

\subsection{Theme A: Usability of Conversation Flow}

Participants liked the idea of conversational interaction, but the flow
was strongest when Proxi either gave clear next steps or when users
already knew the supported phrasing.

\textbf{Recommendation:} Add visible example prompts on the home screen
and after failed requests.

Examples:
\begin{itemize}
    \item ``Try saying: summarize the current result.''
    \item ``Try saying: show my calendar for tomorrow.''
    \item ``Try saying: add this to my notes.''
\end{itemize}

\subsection{Theme B: Usability of System Response}

Users wanted stronger confirmation and recap messages, especially after
multi-step actions.

\textbf{Recommendation:} Improve recap and clarification feedback.

Examples:
\begin{itemize}
    \item ``I found the weather for Hamilton and summarized it below.''
    \item ``I added the meeting and drafted the follow-up email.''
    \item ``I'm not sure which page you mean. Did you mean Project
    Notes?''
\end{itemize}

\subsection{Theme C: Accessibility and Inclusivity}

Elderly and lower-English-fluency users responded positively to Proxi
when the interaction was simple and forgiving.

\textbf{Recommendation:} Add multilingual hints, simplified response
mode, a more visible help area, and an onboarding step for theme and
language preferences.

\subsection{Theme D: Productivity Potential}

Power users and professionals saw the most value when Proxi combined
tasks into one flow.

\textbf{Recommendation:} Promote chained workflows more explicitly.

Examples:
\begin{itemize}
    \item search, summarize, and add to notes.
    \item create meeting and draft email.
    \item read latest email and prepare response.
\end{itemize}

\subsection{Theme E: Exploratory Use and Feature Discovery}

The exploratory free-use session showed that participants were willing
to test Proxi beyond the assigned scenarios, especially when they had
already experienced one or two successful interactions. However, this
also made command discoverability issues more visible. When users were
not given a task prompt, they were more likely to hesitate, test broad
commands, or rely on the help feature.

\textbf{Recommendation:} Improve self-guided exploration by making
feature discovery easier for first-time users.

Examples:
\begin{itemize}
    \item provide a visible set of starter prompts on the home screen,
    \item surface suggested follow-up commands after successful actions,
    \item add a lightweight ``explore'' or ``try asking'' panel,
    \item show capability categories such as email, notes, calendar,
    and weather.
\end{itemize}

\section{Conclusion}

This usability study found that Proxi has strong promise as both an
accessibility tool and a productivity assistant. Participants valued the
ability to complete practical workflows involving email, calendar,
weather, notes, and search through a high-level conversational
interface.

The strongest positive findings were related to reduced application
switching, voice-first convenience, and practical support for everyday
tasks. The most important design challenge was command discoverability:
participants often knew what they wanted to do but were not always sure
how to phrase the request. Better onboarding, stronger examples,
multilingual discoverability, and improved handling of follow-up
references would substantially strengthen the user experience.

The exploratory free-use sessions strengthened these findings by showing
how participants behaved when they were not constrained by scripted
tasks. In that setting, participants naturally explored feature
boundaries, attempted broader requests, and relied more heavily on
their own assumptions about what Proxi could do. This confirmed both
the usefulness of the system and the importance of making feature
discovery more visible for self-directed use.

Overall, the study supports the conclusion that Proxi is functionally
useful, well aligned with its intended audience, and capable of becoming
a highly effective assistant with a few targeted usability improvements.

\appendix

\section{Appendix A: Moderator Script}

\subsection{Introduction}

Hi, thank you for participating in this usability study for Project
Proxi. Today, you will be using Proxi, an AI-powered desktop assistant
that supports voice and text interaction. We are testing the system, not
you. There are no right or wrong answers, and you may stop at any time.

\subsection{Before Starting}

\begin{itemize}
    \item Ask participant to review and sign consent form.
    \item Explain note-taking and optional session recording.
    \item Explain that questions will be asked before, during, and after
    the tasks.
    \item Remind the participant to be honest about confusion,
    frustration, or delight.
\end{itemize}

\subsection{Moderator Prompts During Session}

\begin{itemize}
    \item What are you thinking right now?
    \item Why did you phrase it that way?
    \item What did you expect Proxi to do?
    \item Was that response clear to you?
    \item Would you try that differently now?
    \item Is there anything about this task that feels difficult or
    unnatural?
\end{itemize}

\subsection{Exploratory Session Prompt}

For the next part of the study, you will have open access to Proxi for a
short period of time. There are no required tasks in this part. Please
feel free to try anything that seems useful, interesting, or natural to
you. You can use voice or text. We are interested in what you choose to
do on your own, what you expect Proxi to be able to do, and whether
anything surprises or confuses you.

\section{Appendix B: Full Questionnaires}

\subsection{Post-Test Rating Questions}

Please rate the following on a scale from 1 to 5:
\begin{enumerate}[label=\arabic*.]
    \item Overall, I was satisfied with my experience using Proxi.
    \item I felt that Proxi could understand my requests.
    \item I felt that Proxi was useful.
    \item I would consider using Proxi again.
    \item I would recommend a tool like Proxi to someone who values
    accessibility or efficiency.
\end{enumerate}

\subsection{Final Open-Ended Questions}

\begin{enumerate}[label=\arabic*.]
    \item What did you like most?
    \item What did you dislike most?
    \item What would you improve first?
    \item What feature should Proxi support next?
\end{enumerate}

\section{Appendix C: Data Logging Template}

\begin{table}[H]
\centering
\caption{Sample Data Logging Template}
\begin{tabularx}{\textwidth}{p{1.1cm}p{2.6cm}cccccX}
\toprule
ID & Task & Time & Attempts & Assists & Rating & Notes \\
\midrule
P1 & Search and summarize & 44 s & 1 & 0 & 5 &
Smooth interaction; strong positive value. \\
P1 & Update note page & 58 s & 2 & 0 & 4 &
Minor confusion with page title match. \\
P2 & Weather by voice & 49 s & 1 & 0 & 5 &
Comfortable and direct. \\
P2 & Punjabi follow-up & 65 s & 2 & 1 & 4 &
Positive surprise; wanted earlier visibility. \\
P3 & Create event & 40 s & 1 & 0 & 5 & Natural workflow. \\
P3 & Draft client email & 53 s & 1 & 0 & 5 &
High real-world value. \\
P4 & Weather summary & 46 s & 1 & 0 & 4 & Clear result. \\
P4 & Reminder creation & 57 s & 2 & 0 & 4 &
Small hesitation with phrasing. \\
\bottomrule
\end{tabularx}
\end{table}

\begin{table}[H]
\centering
\caption{Exploratory Session Logging Template}
\begin{tabularx}{\textwidth}{p{1.2cm}p{2.8cm}p{2cm}X}
\toprule
ID & First Feature Tried & Help Used & Notes \\
\midrule
P1 & Search and notes & Yes &
Tested whether Proxi could chain multiple actions. \\
P4 & Weather and reminders & Yes &
Explored daily planning tasks without prompting. \\
\bottomrule
\end{tabularx}
\end{table}

\section{Appendix D: Sample Quote Bank}

\subsection{Raj}

\begin{itemize}
    \item ``This feels useful because it cuts out the app switching.''
    \item ``The note update was good, but it should be more forgiving
    with names.''
    \item ``I can see myself using this for quick research summaries.''
\end{itemize}

\subsection{Tony}

\begin{itemize}
    \item ``I like the light one better. It feels simpler.''
    \item ``Oh, I didn't know Punjabi option was also there.''
    \item ``If it understands both, that is very good.''
    \item ``Typing is more difficult for me than just saying it.''
    \item ``I want it to explain slowly when I ask again.''
\end{itemize}

\subsection{Jassi}

\begin{itemize}
    \item ``This is actually practical for my work.''
    \item ``The calendar part felt the most natural.''
    \item ``I liked that it asked before doing something important.''
    \item ``If I'm driving between showings, this would help.''
\end{itemize}

\subsection{Mikayla}

\begin{itemize}
    \item ``Weather plus reminder is a nice combo.''
    \item ``I didn't know what exact words to use at first.''
    \item ``Help made sense, but I wish I saw examples earlier.''
    \item ``It felt smart when it remembered the earlier thing.''
\end{itemize}

\end{document}